% Invented text. A representative deck: a title frame, stepped bullets,
% nested lists, a table, columns, a formula, an image, a code block whose
% longest line does not fit, and a frame with more prose than a slide holds
% — the frames a talk sets.
\documentclass[aspectratio=169]{beamer}
\usepackage{booktabs}
\usepackage{graphicx}
\fonts{ dir = "../corpus/fonts", sans = "Fira Sans", mono = "Source Code Pro",
  math = "Fira Math" }

\title{Bread at Altitude}
\subtitle{Notes from a village bakery}
\author{A. Placeholder}
\date{}

\begin{document}

\maketitle

\begin{frame}{Why altitude matters}
  \begin{itemize}
    \item<1-> Water boils at a lower temperature, so steam forms sooner.
    \item<2-> Dough rises faster in thinner air, then collapses.
    \item<3-> Flour dries out: the same recipe needs more water.
  \end{itemize}
\end{frame}

\begin{frame}{Adjusting a recipe}
  \begin{enumerate}
    \item Reduce the yeast by a quarter.
    \item Add water a spoonful at a time:
    \begin{itemize}
      \item until the dough is tacky, not sticky;
      \item more in winter, when the air is driest.
    \end{itemize}
    \item Bake hotter and shorter.
  \end{enumerate}
\end{frame}

\begin{frame}{Three loaves compared}
  \centering
  \begin{tabular}{@{}lrrl@{}}
    \toprule
    \textbf{Loaf} & \textbf{Rise (min)} & \textbf{Bake (min)} & \textbf{Crumb} \\
    \midrule
    Plain white & 45 & 30 & open, even \\
    Rye blend & 70 & 40 & dense, moist \\
    Seeded wholemeal & 60 & 35 & tight, nutty \\
    \bottomrule
  \end{tabular}
\end{frame}

\begin{frame}{Two methods side by side}
  \begin{columns}
    \begin{column}{0.48\textwidth}
      \textbf{Straight dough}
      \begin{itemize}
        \item one mixing, one rise
        \item quick, but less flavour
      \end{itemize}
    \end{column}
    \begin{column}{0.48\textwidth}
      \textbf{Sponge and dough}
      \begin{itemize}
        \item a starter rests overnight
        \item slower, with a deeper flavour
      \end{itemize}
    \end{column}
  \end{columns}
\end{frame}

\begin{frame}{How much water}
  The water a dough needs grows with the dryness of the flour:
  \[
    w = w_0 \left( 1 + \frac{h_0 - h}{100} \right)
  \]
  where $w_0$ is the recipe's water, $h_0$ the humidity it was written
  for, and $h$ the humidity in the bakery today.
\end{frame}

\begin{frame}{The proving shelf}
  \centering
  \includegraphics[width=0.4\textwidth, alt={Three coloured rectangles for
  three proving trays}]{../corpus/rects.png}
\end{frame}

\begin{frame}[fragile]{The oven timer}
\begin{verbatim}
def bake_minutes(base, altitude_m):
    return base - altitude_m // 500 if altitude_m > 1000 else base  # shorter bakes high up, longer bakes in the valley
\end{verbatim}
\end{frame}

\begin{frame}{What the bakers said}
  Every baker interviewed had learnt the adjustments by trial rather than
  from a book, and most had kept notes of their failures longer than of
  their successes. The oldest kept a ledger of every batch for thirty years,
  with the weather, the flour and the result written beside each one. Read
  in order, the ledger shows the recipes drifting slowly towards less yeast
  and more water, a season at a time, until they settled. Nobody had planned
  the change; it accumulated from one small correction after another,
  each made because yesterday's loaf had been a little too flat or a little
  too dry, and each kept because the next loaf was better.
\end{frame}

\end{document}
